% ============================================================================
% Chapter 1 — Introduction
% ============================================================================
\section{Introduction}
\label{sec:introduction}

Senior leaders routinely report that the years in which their formal
authority, reputation, and network---their \emph{career capital}---peaked were
also the years in which their capacity to do the work that built that capital
collapsed. Surveys of executives find majorities considering exits for
well-being reasons \citep{deloitte2022wellbeing}, median CEO tenures have
shortened \citep{conference2023tenure}, and managerial burnout has become a
recognized organizational risk rather than a private misfortune
\citep{maslach2001job,veldhuis2020burnout,muller2026systematic,tu2025dancing}.
Existing accounts treat this pattern either as a resource problem (demands
exceed resources; \citealp{demerouti2001job,bakker2017jdr}), as a stress
process \citep{hobfoll1989conservation}, or as an organizational-level
capability trap \citep{repenning2001nobody,repenning2002capability}. What is
missing is a formal, individual-level theory of \emph{why success itself}
should systematically produce the depletion---and of which interventions can
and cannot reverse it.

This paper supplies that theory. Dynamic Leadership Vitality Theory (DLVT) is
the first closed-form, individual-level dynamical model in which accumulated
career capital endogenously generates the organizational complexity that
drains deployment vitality, yielding a \emph{globally} asymptotically stable
low-vitality/high-status attractor---and, under a power-law load kernel, an
intervention asymmetry in which equilibrium vitality is invariant to the
scope-coupling parameter $\beta$ and movable only through recovery-side
parameters.

The model is deliberately minimal: two coupled ordinary differential
equations. Vitality $V(t)$ recovers toward a ceiling $\Vmax$ and is drained at
a rate that grows nonlinearly with organizational complexity $O(t)$; career
capital $C(t)$ accumulates through energy-gated impact and depreciates; and
complexity is an increasing function of the leader's own capital,
$O = O_0 + \beta C^{\eta}$. The loop closes: capital builds complexity,
complexity drains vitality, vitality gates the impact that builds capital.
The system settles---from essentially every initial condition---into a single
interior equilibrium at which status is high and vitality is chronically low.

\paragraph{Position relative to capability-trap models.}
DLVT is closest in spirit to the system-dynamics literature on capability
traps and capability erosion
\citep{repenning2002capability,repenning2002simulation,rahmandad2016capability,sterman2000business}.
Those models are organizational-level and, crucially, \emph{bistable}: firms
tip into the trap when short-run pressure crosses a threshold, and a
sufficiently large push can tip them back out. DLVT's contrast is sharp. At
the individual level, with the mechanisms modeled here, the low-vitality
equilibrium is not one basin of two---it is the unique attractor of the
positive orthant. The trap is the \emph{default destination}, not an
unfortunate branch. This difference is not rhetorical: it is a falsifiable
structural prediction (no multistability, no hysteresis, no tipping point in
the baseline kernel family; \S\ref{sec:propositions}), and it reverses the
managerial intuition that a one-time ``reset'' (a sabbatical, a
reorganization) can durably relocate the system.

\paragraph{Micro-foundations.}
The vitality construct and its dynamics are grounded in the effort--recovery
model of work strain \citep{meijman1998effort} and the literature on
subjective vitality as a measurable, fluctuating energetic resource
\citep{ryan1997subjective,nix1999revitalization,quinn2012humanenergy}. We
explicitly do \emph{not} build on ego-depletion theory
\citep{baumeister1998ego}: the depletion effect has failed meta-analytic
scrutiny and multisite preregistered replication
\citep{hagger2010ego,carter2015series,vohs2021multisite}, and a theory whose
core state variable rested on it would inherit that fragility. The
effort--recovery mechanism---load during effort, restoration during respite,
and cumulative strain when recovery is structurally insufficient---requires no
depletable willpower substance and is the empirically surviving foundation.

\paragraph{Relation to human-capital theory.}
DLVT does not compete with \citet{becker1993human}; it \emph{adds} to that
framework an energetic-deployment constraint that human-capital theory does
not model. In Becker's account, accumulated capital is deployed at will; in
DLVT, deployment requires vitality, and the accumulation process itself
erodes the vitality needed for deployment. The result is a wedge between
capital owned and capital usable---formalized here as the energy-gated impact
function---that grows with seniority.

\paragraph{Formalism.}
Methodologically, the paper follows the tradition of small, fully analyzed
dynamical models in adjacent fields: ecology
\citep{may1976simple,murray2002mathematical,scheffer2001catastrophic},
organizational simulation \citep{pluchino2010peter}, and computational
psychopathology \citep{cramer2016depression}. Every analytical claim below is
accompanied by a proof or proof sketch, and every numerical claim is pinned by
an automated test in the companion code, so that manuscript and implementation
cannot silently drift apart.

\paragraph{A note on terminology.}
Earlier drafts of this work labeled the low-vitality/high-status equilibrium
the ``zombie-leader equilibrium.'' We retire ``zombie'' as a defined
theoretical construct: it imported connotations the mathematics does not
support and invited misreading of the equilibrium as a discrete type of
person rather than a region of a continuous state space. Throughout this
corrected manuscript the object is called the \emph{low-vitality equilibrium}
(Definition~\ref{def:lowvitality}); the informal label survives only in
legacy function names in the companion code (e.g., \texttt{is\_zombie}).

\paragraph{Contributions.}
The paper makes four contributions.
(i)~\emph{Theory:} a closed-form, individual-level capability-trap model with
a single globally stable attractor, positioned against the bistable
organizational-trap tradition (\S\ref{sec:foundations},
\S\ref{sec:formulation}).
(ii)~\emph{Corrected mathematics:} existence and uniqueness of the interior
equilibrium in the entire positive orthant (Theorem~\ref{thm:uniqueness}),
trace negativity ruling out Hopf bifurcations (Lemma~\ref{lem:trace}), and
global asymptotic stability on $\{C>0\}$ via a corrected trapping rectangle
(Theorem~\ref{thm:gas}); several claims of earlier drafts are explicitly
retracted and replaced (\S\ref{sec:transient}, Appendices~\ref{app:bifurcation}
and \ref{app:gas}).
(iii)~\emph{Intervention asymmetry:} under the power-law load kernel,
equilibrium vitality is invariant to the scope coupling $\beta$
(Corollary~\ref{cor:scope}); scope reduction alone cannot raise equilibrium
vitality, while recovery-side parameters can---stated honestly as a property
of the kernel family, not a universal law.
(iv)~\emph{Falsifiability:} six propositions with declared scope conditions
and an empirical roadmap culminating in a pre-registered $2{\times}2$
scope-versus-recovery field test (\S\ref{sec:propositions},
\S\ref{sec:discussion}).

\paragraph{Roadmap.}
Section~\ref{sec:foundations} develops the theoretical foundations.
Section~\ref{sec:formulation} states the model, the baseline calibration
(Table~\ref{tab:parameters}), and the formal results.
Section~\ref{sec:results} reports the numerical analysis and the proposition
set. Section~\ref{sec:discussion} discusses limitations---including the
illustrative status of the calibration and the stipulated strategic
threshold---and the validation roadmap. Appendices~A1 and A4--A10 contain the
complete proofs and numerical audits.
